\documentclass[../../../include/open-logic-section]{subfiles}

\begin{document}

\olfileid{sth}{ordinals}{replacement}
\olsection{تعويض}

په \olref[sth][ordinals][vn]{sec} کښې مو د ترتيبي عددونو د معرفي کولو
وجه دا وړاندې کړه چې د ترتيب د ډولونو په توګه، يعنې د ښه ترتيبونو د
کانوني استازو په توګه، ئې کارولے شو. د دې د عملي کېدو دپاره به ثابتول
پکار وي چې \emph{هر ښه ترتيب له کوم ترتيبي عدد سره آيزومورف دے}.
په دې سره $\ordtype{A, <}$ د هغه ترتيبي عدد $\alpha$ په توګه تعريفولے
شو چې $\tuple{A, <} \isomorphic \alpha$ وي.

له بده مرغه، يوازې د هغو بديهي اصلونو په کارولو چې تر اوسه مو معرفي
کړي دي، مطلوبه نتيجه \emph{نۀ شو} ثابتولے. (وجه به ئې په
\olref[replacement][strength]{sec} کښې ووينو؛ خو اوس مطلب دا دے چې
نۀ ئې شو ثابتولے.) نوے فکر پکار دے، او هغه دا دے:

\begin{axiom}[د تعويض شېما]
د هر فارمول $\phi(x, y)$ دپاره، لاندې خبره بديهي اصل دے:
\begin{quote}
	د هر $A$ دپاره، که $(\forall x \in A)\lexists![y][\phi(x,y)]$ وي، نو
	$\Setabs{y}{(\exists x \in A)\phi(x,y)}$ شتون لري.
\end{quote}
\end{axiom}
\noindent
د بېلولو په شان، دا يوه شېما ده: د بې شمېره بېلو $\phi$ ګانو د هر يوه
دپاره يو بديهي اصل ورکوي، نو نامتناهي ډېر بديهي اصلونه ترې جوړېږي.
په هماغه معنا داسې هم ليکل کېدلے شي، او عموماً همداسې ليکل کېږي:

\begin{defish}
د هر داسې فارمول $\phi(x,y)$ دپاره چې «$B$» پکښې نۀ راځي، لاندې خبره بديهي اصل دے:
\[
\forall A[(\forall x \in A)\lexists![y][\phi(x,y)] \lif \exists B\forall y (y \in B \liff (\exists x \in A)\phi(x,y))]
\]
\end{defish}

خو په لومړي ځل دا فارمول څه پېچلے ښکاري. د تعويض لاندې لنډه پايله ښايي
زموږ مطلوب شهودي تصور په \emph{لا روښانه} توګه بيان کړي:\footnote{ترم هغه
عبارت دے چې د خپلو ازادو متغيرونو په هر بشپړ ټاکلو کښې عين يو شے ټاکي.
د بېلګې په توګه، «$\emptyset$» هغه ترم دے چې تش سټ ټاکي؛ «$\{x\}$» هغه
ترم دے چې د $x$ يوغړيز سټ ټاکي، هر څه چې $x$ وي؛ او «$x \cup y$» هغه
ترم دے چې د $x$ او $y$ اتحاد ټاکي، هر څه چې دا دواړه وي.}

\begin{cor}
د هر ترم $\tau(x)$ او هر سټ $A$ دپاره، دا سټ شتون لري:
\[
	\Setabs{\tau(x)}{x \in A} = \Setabs{y}{(\exists x \in A)y = \tau(x)}.
\]
\end{cor}

\begin{proof}
ځکه چې $\tau$ يو \emph{ترم} دے، $\forall x \lexists![y][\tau(x) = y]$ دے.
په تېره بيا، $(\forall x \in A)\lexists![y][\tau(x) = y]$ دے. نو
$\Setabs{y}{(\exists x \in A)\tau(x) = y}$ د تعويض له مخې شتون لري.
\end{proof}
\noindent
له دې ښکاري چې «تعويض» د دې بديهي اصل دپاره ښه نوم دے: که سټ $A$ درکړے
شي، نو د $\Setabs{\tau(x)}{x \in A}$ په توګه نوے سټ جوړولے شئ، په دې
چې د $A$ هر غړے د~$\tau$ له مخې د هغه په تصوير بدل کړئ. په رښتيا، د
تابعې له مخې د سټ د تصوير د نښو په تعقيب،
$\funimage{\tau}{A}$ د $\Setabs{\tau(x)}{x \in A}$ دپاره ليکلے شو.

خو دا ډېره مهمه ده چې $\tau$ يو \emph{ترم} دے. ضرور نۀ ده چې د
\olref[sfr][fun][rel]{sec} په معنا د کومې \emph{تابعې} نوم وي، يعنې د
مرتبو جوړو د کوم ټاکلي سټ نوم. ځکه که $f$ په دې معنا تابعه وي، نو سټ
$\funimage{f}{A} = \Setabs{f(x)}{x \in
A}$ يوازې د $\ran{f}$ يو ټاکلے فرعي سټ دے، او د~$\Zminus$ بديهي
اصلونه لا له وړاندې ئې شتون تضمينوي.\footnote{يوازې
$\Setabs{y \in \bigcup \bigcup f}{(\exists x \in A)y =
f(x)}$ په پام کښې ونيسئ.} خو تعويض زموږ بديهي اصلونو ته يوه
\emph{قوي} اضافه ده، لکه چې په \olref[sth][replacement][]{chap} کښې به ووينو.

\end{document}
